\section*{Original Section 3.3.3: Derandomizing Noether's Normalization Lemma (journal pp.37--39)}
We now prove Theorem 1.12 and Theorem 1.13. For convenience, we restate the theorems below.
\begin{result}{Theorem 1.12}
Let $n,d\in\mathbb N^+$, $r\in\{0,1,\ldots,n\}$, $k=n-r-1$, and $\epsilon\in(0,1)$. There exists an explicit collection $\mathcal L$ of linear maps $\Aspace^{n+1}\to\Aspace^{r+1}$ of size $\operatorname{poly}(N(k,d,n),n,1/\epsilon)$ such that for every subvariety $\mathcal V\subseteq\PP^n$ of dimension $r$ and degree at most $d$, all but at most $\epsilon$-fraction of $\pi\in\mathcal L$ induce a surjective finite morphism from $\mathcal V$ to $\PP^r$.\footnote{Let $N(k,d,n)=1$ when $r=n$ (i.e., $k=-1$). Similarly, in Theorem 1.13, let $N(k,d,n-1)=1$ when $r=n$ or $r=0$ (i.e., $k=-1$ or $k=n-1$).} Moreover, $\mathcal L$ can be computed in time polynomial in $|\mathcal L|$ (and $\log p$, if $\operatorname{char}(\FF)=p>0$).
\end{result}
\begin{proof}
If $r=n$, just use the identity map on $\Aspace^{n+1}=\Aspace^{r+1}$. So assume $r<n$. Use Theorem 3.14 to construct an $(n,d,\epsilon)$-evasive $k$-subspace family $\mathcal H$ on $\PP^n$ of size $\operatorname{poly}(N(k,d,n),n,1/\epsilon)$. For each $k$-subspace $W\in\mathcal H$, compute a surjective linear map $\pi_W:\Aspace^{n+1}\to\Aspace^{r+1}$ such that $W=\PP\ker(\pi_W)$. Let $\mathcal L=\{\pi_W:W\in\mathcal H\}$. Then $\mathcal L$ is a desired collection of linear maps by Lemma 3.12.
\end{proof}
\begin{result}{Theorem 1.13}
Let $n,d\in\mathbb N^+$, $r\in\{0,1,\ldots,n\}$, $k=n-r-1$, and $\epsilon\in(0,1)$. There exists an explicit collection $\mathcal L$ of linear maps $\Aspace^n\to\Aspace^r$ of size $\operatorname{poly}(N(k,d,n-1),n,1/\epsilon)$ such that for every subvariety $\mathcal V\subseteq\Aspace^n$ of dimension $r$ and degree at most $d$, all but at most $\epsilon$-fraction of $\pi\in\mathcal L$ restrict to a surjective finite morphism from $\mathcal V$ to $\Aspace^r$. Moreover, $\mathcal L$ can be computed in time polynomial in $|\mathcal L|$ (and $\log p$, if $\operatorname{char}(\FF)=p>0$).
\end{result}
\begin{proof}
If $r=n$, just use the identity map on $\Aspace^n=\Aspace^r$. If $r=0$, use the only map $\Aspace^n\to\Aspace^0$. So assume $0<r<n$. Regard $\Aspace^n$ as an open subset of $\PP^n$ via $(x_1,\ldots,x_n)\mapsto(x_1,\ldots,x_n,1)$. Similarly, regard $\Aspace^r$ as an open subset of $\PP^r$ via $(x_1,\ldots,x_r)\mapsto(x_1,\ldots,x_r,1)$. Let $H_\infty$ be the hyperplane at infinity of $\PP^n$ defined by $X_{n+1}=0$.
Use Theorem 3.14 to construct an $(n-1,d,\epsilon)$-evasive $k$-subspace family $\mathcal H$ on $H_\infty\cong\PP^{n-1}$ of size $\operatorname{poly}(N(k,d,n-1),n,1/\epsilon)$. For each $W\in\mathcal H$, choose $n-k=r+1$ homogeneous linear polynomials $\ell_1,\ldots,\ell_{r+1}\in\FF[X_1,\ldots,X_{n+1}]$ such that $\ell_{r+1}=X_{n+1}$, $\ell_1,\ldots,\ell_r\in\FF[X_1,\ldots,X_n]$, and $W=V(\ell_1,\ldots,\ell_{r+1})$. This is possible as $W\subseteq H_\infty=V(X_{n+1})$.
These $r+1$ linear polynomials determine a projection $\pi_W:\PP^n\setminus W\to\PP^r$, defined by
\begin{align*}
x=(x_1,\ldots,x_{n+1})&\mapsto(\ell_1(x),\ldots,\ell_{r+1}(x))\\
&=(\ell_1(x),\ldots,\ell_r(x),x_{n+1}).
\end{align*}
As $x_{n+1}=1$ for $x\in\Aspace^n$, we have $\pi_W(\Aspace^n)\subseteq\Aspace^r$. Restricting $\pi_W$ on $\Aspace^n$ yields a map $\pi_W|_{\Aspace^n}:\Aspace^n\to\Aspace^r$, which is a linear map as $\ell_1,\ldots,\ell_r$ are homogeneous linear polynomials in $\FF[X_1,\ldots,X_n]$. Let $\mathcal L=\{\pi_W|_{\Aspace^n}:W\in\mathcal H\}$.
Let $\mathcal V$ be a subvariety of $\Aspace^n$ of dimension $r$ and degree at most $d$. Its projective closure $\mathcal V^{\rm cl}$ has dimension $\dim(\mathcal V)=r$ and degree $\deg(\mathcal V)\le d$. By the definition of $\mathcal V^{\rm cl}$, none of the irreducible components of $\mathcal V^{\rm cl}$ is fully contained in $H_\infty$. So by Lemma 2.7, the projective subvariety $\mathcal V^{\rm cl}\cap H_\infty$ has dimension $r-1$ and degree at most $d$.
By the choice of $\mathcal H$, all but at most $\epsilon$-fraction of $W\in\mathcal H$ are disjoint from $\mathcal V^{\rm cl}\cap H_\infty$ and hence from $\mathcal V^{\rm cl}$. So we just need to prove that for every $W\in\mathcal H$ disjoint from $\mathcal V^{\rm cl}$ and the corresponding projection $\pi_W$, the map $\pi_W|_{\mathcal V}:\mathcal V\to\Aspace^r$ is a surjective finite morphism. This follows from Lemma 3.12 and the fact that $\mathcal V=\mathcal V^{\rm cl}\cap(\pi_W)^{-1}(\Aspace^r)$.
\end{proof}
